% Part: lambda-calculus
% Chapter: lambda-definability
% Section: primitive-recursive-functions

\documentclass[../../../include/open-logic-section]{subfiles}

\begin{document}

\olfileid{lam}{ldf}{prf}
\olsection{Fungsi Rekursif Primitif iku \usetoken{S}{lambda definable}}

Elinga yen fungsi rekursif primitif yaiku fungsi-fungsi kang bisa
ditetepake saka fungsi dhasar $\Zero$, $\Succ$, lan $\Proj{n}{i}$
kanthi komposisi lan rekursi primitif.

\begin{lem}
  \ollabel{lem:basic}
  Fungsi rekursif primitif dhasar $\Zero$, $\Succ$, lan
  proyeksi~$\Proj{n}{i}$ iku !!{lambda definable}.
\end{lem}

\begin{proof}
Fungsi-fungsi mau !!{lambda define}d dening suku-suku iki:
\begin{align*}
  \fn{Zero} & \ident \lambd[a][\lambd[fx][x]]\\
  \fn{Succ} & \ident \lambd[a][\lambd[fx][f (a f x)]] \\
  \fn{Proj}^n_i & \ident \lambd[x_0\dots x_{n-1}][x_i]
\end{align*}
\end{proof}

\begin{lem}
  \ollabel{lem:comp} Upamakna fungsi kanthi aritas $k$, yaiku $f$,
  lan fungsi-fungsi kanthi aritas $n$, yaiku $g_0, \dots, g_{k-1}$,
  !!{lambda definable} dening suku $F$, $G_0$, \dots, $G_{k-1}$,
  lan $h$ ditetepake saka fungsi-fungsi mau kanthi komposisi.
  Mula $h$ uga !!{lambda definable}.
\end{lem}

\begin{proof}
  $h$ bisa !!{lambda define}d dening suku
  \[
  H \ident \lambd[x_0 \dots x_{n-1}][F\, (G_0 x_0 \dots
    x_{n-1}) \dots (G_{k-1} x_0 \dots x_{n-1})]
  \]
  Verifikasi pratelan iki kita tinggalake minangka latihan.
\end{proof}

\begin{prob}
  Rampungna bukti \olref[lam][ldf][prf]{lem:comp} kanthi nuduhake
  yen $H\num{n_0}\dots\num{n_{n-1}} \red \num{h(n_0, \dots,
    n_{n-1})}$.
\end{prob}

Gatekna yen \olref{lem:comp} ora mbutuhake $f$ lan $g_0$,
\dots,~$g_{k-1}$ minangka fungsi rekursif primitif; kang dibutuhake
mung fungsi-fungsi iku total lan !!{lambda definable}.

\begin{lem}\ollabel{lem:prim}
  Upamakna $f$ fungsi kanthi aritas $n$ lan~$g$ fungsi kanthi
  aritas $n+2$. Loro-lorone !!{lambda definable} dening suku $F$
  lan~$G$, dene fungsi~$h$ ditetepake saka $f$ lan $g$ kanthi
  rekursi primitif. Mula $h$ uga !!{lambda definable}.
\end{lem}

\begin{proof}
  Elinga yen $h$ ditetepake dening
  \begin{align*}
    h(x_1, \dots, x_n, 0) &= f(x_1, \dots, x_n)\\
    h(x_1, \dots, x_n, y+1) & = g(x_1, \dots, x_n, y, h(x_1, \dots, x_n, y)).
  \end{align*}
  Kanthi informal, definisi rekursif primitif ngetrapake fungsi
  pambaru kahanan saka $g$ ping $y$, kawiwitan saka $f(x_1,
  \dots, x_n)$. Iki mirip definisi numeral Church, kang uga
  nggunakake iterasi.

  Supaya luwih prasaja, kita menehi definisi lan bukti kanggo siji
  argumen tambahan~$x$. Fungsi $h$ !!{lambda define}d dening:
  \begin{align*}
    H \ident & \lambd[x][\lambd[y][\fn{Snd} (y
        D \tuple{\num{0}, F x})]]
    \intertext{kanthi }
    D \ident & \lambd[p][\tuple{\fn{Succ} (\fn{Fst}\, p), (G x
          (\fn{Fst}\, p) (\fn{Snd}\, p))}]
  \end{align*}
  Kahanan iterasi kang kita simpen iku pasangan: komponen kapisan
  nuduhake nilai $y$ saiki, dene komponen kapindho nuduhake nilai
  $h$ kang cocog. Ing saben langkah iterasi, kita gawe pasangan
  anyar kanggo nilai $y$ lan $h$; sawuse iterasi rampung, komponen
  kapindho dibalekake minangka nilai pungkasan $h$. Suku pambaru
  dimaksudake minangka cithakan kang peubah paramètere bebas lan
  kaiket nalika cithakan iki dilebokake ing suku utama. Saiki kita
  buktekake yen suku iki pancen makili rekursi primitif.

  Kita arep mbuktekake yen kanggo saben $n$ lan $m$, $H\,\num{n}\,\num{m} \red
  \num{h(n,m)}$. Dhisik kita tuduhake yen yen $D_n \ident
  \Subst{D}{\num{n}}{x}$, mula $D_n^m \tuple{\num{0}, F\, \num{n}} \red
  \tuple{\num{m}, \num{h(n, m)}}$. Kita nggunakake induksi marang~$m$.

  Yen $m=0$, kang kudu dituduhake yaiku $D_n^0 \tuple{\num{0}, F\, \num{n}} \red
  \tuple{\num{0}, \num{h(n, 0)}}$. Nanging $D_n^0 \tuple{\num{0}, F\,
    \num{n}}$ pancen mung $\tuple{\num{0}, F\, \num{n}}$. Amarga $F$
  !!{lambda define}s~$f$, suku iki direduksi dadi $\tuple{\num{0},
    \num{f(n)}}$; lan amarga $f(n) = h(n, 0)$, asile $\tuple{\num{0},
    \num{h(n,0)}}$.

  Saiki upamakna $D_n^m \tuple{\num{0}, F\, \num{n}} \red
  \tuple{\num{m}, \num{h(n, m)}}$. Kita arep
  nuduhake yen $D_n^{m+1} \tuple{\num{0}, F\, \num{n}} \red
  \tuple{\num{m+1}, \num{h(n, m+1)}}$.
  \begin{align*}
    D_n^{m+1} \tuple{\num{0}, F\, \num{n}}
    & \ident D_n(D_n^{m} \tuple{\num{0}, F\, \num{n}})\\
    & \red D_n\,\tuple{\num{m}, \num{h(n, m)}} \text{\qquad (miturut HI)}\\
    & \ident (\lambd[p][
      \tuple{
        \fn{Succ} (\fn{Fst}\, p), (G \, \num{n} (\fn{Fst}\, p) (\fn{Snd}\, p))
    }]) \tuple{\num{m}, \num{h(n,m)}}\\
    & \redone \tuple{\fn{Succ} (\fn{Fst}\, \tuple{\num{m}, \num{h(n,m)}}),
      (G \, \num{n} (\fn{Fst}\, \tuple{\num{m}, \num{h(n,m)}}) (\fn{Snd}\, \tuple{\num{m}, \num{h(n,m)}}))}\\
    & \red 
    \tuple{\fn{Succ}\, \num{m}, (G \, \num{n} \,\num{m} \, \num{h(n,m)})}\\
    & \red \tuple{\num{m+1}, \num{g(n, m, h(n, m))}}
  \end{align*}
  Amarga $g(n, m, h(n, m)) = h(n, m+1)$, bukti rampung.

  Pungkasan, gatekna
  \begin{align*}
    H\,\num n\,\num m &\ident \lambd[x][\lambd[y][\fn{Snd} (y
        (\lambd[p][\tuple{\fn{Succ} (\fn{Fst}\, p), (G\, x\,
          (\fn{Fst}\, p)\, (\fn{Snd}\, p))}]) \tuple{\num{0}, F x})]]\\
    & \qquad\,\num n\, \num m\\
    & \red \fn{Snd} (\num m \,
    \underbrace{(\lambd[p][\tuple{\fn{Succ} (\fn{Fst}\, p), (G \,\num n\,
      (\fn{Fst}\, p) (\fn{Snd}\, p))}])}_{D_n} \tuple{\num{0}, F \num n})\\
    & \ident \fn{Snd} (\num m \, D_n\, \tuple{\num{0}, F \num n})\\
    & \red \fn{Snd} \,(D_n^m  \tuple{\num{0}, F \num n})\\
    & \red \fn{Snd} \,\tuple{\num{m}, \num{h(n, m)}} \\
    & \red \num{h(n,m)}. 
  \end{align*}  
\end{proof}


\begin{prop}
  Saben fungsi rekursif primitif iku !!{lambda definable}.
\end{prop}

\begin{proof}
  Miturut \olref{lem:basic}, kabeh fungsi dhasar iku !!{lambda definable},
  lan miturut \olref{lem:comp} lan \olref{lem:prim}, kelas fungsi kang
  !!{lambda definable} katutup marang komposisi lan rekursi primitif.
\end{proof}

\end{document}
